% META: {"id": "TEXP-ARI-CO-024", "chapitre": "TEXP-ARITHMETIQUE", "type_objet": "corrige", "exercice_ref": "TEXP-ARI-EX-024", "capacites_codes": ["C7"], "sources_inspiration": [], "status": "generated"}
% BEGIN-VERIFY
% import math
% from sympy import factorint, divisor_count
% assert factorint(360) == {2: 3, 3: 2, 5: 1}
% assert factorint(2520) == {2: 3, 3: 2, 5: 1, 7: 1}
% assert divisor_count(360) == 24
% assert divisor_count(2520) == 48
% assert math.gcd(360, 2520) == 360
% assert 2520 % 360 == 0
% END-VERIFY
\begin{corrige}{TEXP-ARI-EX-024}
\textbf{1.} $360=2^3\times3^2\times5$ et $2\,520=2^3\times3^2\times5\times7$.

\textbf{2.} Le nombre de diviseurs est le produit des exposants augmentés de
$1$ : pour $360$, $(3+1)(2+1)(1+1)=24$ ; pour $2\,520$,
$(3+1)(2+1)(1+1)(1+1)=48$.

\textbf{3.} Le PGCD s'obtient en prenant, pour chaque premier, le plus petit
exposant : $2^3\times3^2\times5=360$. Comme $360$ divise $2\,520$, c'est
cohérent.
\end{corrige}
